\par\clearpage
\noindent\makebox[0pt][l]{\fontsize{13}{16}\selectfont 448}\hfill{\small\textbf{POSTHUMOUS PAPERS.}}\hfill\null\par\medskip
\begingroup
\setlength{\abovedisplayskip}{5pt}\setlength{\belowdisplayskip}{5pt}
\[
v-\gamma=a+b,
\]
\noindent or, setting \(\gamma=\mu v\),
\[
v(1-\mu)=a+b.
\]
\par\noindent\hspace*{\parindent}Since
\[
Rv=r
\]
\noindent holds, it follows that:]
\[
(1-\mu)r=R(a+b).
\]
\par\noindent\hspace*{\parindent}[Moreover, to a close approximation,
\[
\begin{aligned}
A-B&=(2R+A+B)\operatorname{tang}\tfrac12(a-b)\mathbin{.}\operatorname{tang}\tfrac12v,\\
\tfrac{A-B}{2R}&=\tfrac12(a-b)\mathbin{.}\tfrac12\tfrac rR
\end{aligned}
\]
\noindent or]
\[
(a-b)r=2(A-B).
\]
\par\noindent\hspace*{\parindent}Thus if \(b=0\) and \(B=0\), [it follows that]
\[
\begin{aligned}
(1-\mu)r&=Ra\\
ar&=2A;\\
A&=\tfrac{aaR}{2(1-\mu)},\\
(1-\mu)rr&=2AR.
\end{aligned}
\]
\[
\begin{array}{@{}r@{\ }l@{\quad}l@{\quad}r@{\ }l@{\ }l@{}}
\text{[For }\mu=] & 0{,}12 & \text{[we have]} & r= & 3804\,\surd A & \text{[metres]}\\
 &0{,}13&&&3826\,\surd A&\\
 &0{,}14&&&3848\,\surd A&\\
 &0{,}15&&&3870\,\surd A&\\
 &0{,}16&&&3893\,\surd A.&
\end{array}
\]
\par\medskip
\begin{center}
[4.]\par\smallskip
[Formula for calculating heights.]
\end{center}
\par\noindent\hspace*{\parindent}To calculate the heights as accurately as possible, the following formula must be used:
\[
h'-h=\tfrac12(a'-a)\,{}^3\!\surd(\sec a\sec a')\mathbin{.}\{r+\tfrac12(h+h')c\}.
\]
\par\noindent\hspace*{\parindent}Here \(a\) and \(a'\) denote the angles of elevation; \(h\) and \(h'\) the heights above
sea level; \(r\) the distance reduced to sea level; \(c\) the curvature.
\endgroup

\par\clearpage
\noindent\null\hfill{\small\textbf{HEIGHT MEASUREMENTS.}}\hfill\makebox[0pt][r]{\fontsize{13}{16}\selectfont 449}\par\medskip
\begin{center}
[5.]\par\smallskip
[Adjustment of the heights of the principal triangulation points.]
\end{center}
\par\noindent\hspace*{\parindent}[The height measurements were carried out between 1821 and 1825 in the course of the
angle measurements for the arc measurement, using a \textsc{Borda} circle.
They give the height differences between the endpoints of 71 triangle sides
(no height measurements were made along the sides Bremen-Bottel and Hamburg-Lüneburg or at or towards the triangulation point
Inselsberg), from which adjustment is to yield
the height differences between 32 triangulation points.
\par\noindent\hspace*{\parindent}For this purpose the following approximate values are initially assumed
for the heights of the triangulation points.]
\par\medskip
{\centering
\setlength{\tabcolsep}{.85em}
\renewcommand{\arraystretch}{1.1}
\begin{tabular}{|r|l|r||r|l|r|}
\noalign{\hrule height .8pt}
\multicolumn{1}{|c|}{\small No.}&\multicolumn{1}{c|}{\small Station}&
\multicolumn{1}{c||}{\shortstack{\rule{0pt}{1pt}\\\small Assumed\\\small height\\[3pt]\small (metres)}}&
\multicolumn{1}{c|}{\small No.}&\multicolumn{1}{c|}{\small Station}&
\multicolumn{1}{c|}{\shortstack{\rule{0pt}{1pt}\\\small Assumed\\\small height\\[3pt]\small (metres)}}\\[2pt]
\hline
1 & Langwarden & 24,726 & 17 & Lüneburg & 76,832\\
2 & Jever & 33,419 & 18 & Nindorf & 116,707\\
3 & Varel & 38,776 & 19 & Timpenberg & 117,140\\
4 & Bremerlehe & 25,756 & 20 & Wulfsode & 104,510\\
5 & Garlste & 50,590 & 21 & Hauselberg & 120,443\\
6 & Brillit & 44,584 & 22 & Breithorn & 120,499\\
7 & Bremen & 82,882 & 23 & Falkenberg & 150,805\\
8 & Zeven & 44,820 & 24 & Scharnhorst & 93,957\\
9 & Steinberg & 72,336 & 25 & Garssen & 77,079\\
10 & Brüttendorf & 49,631 & 26 & Deister & 307,912\\
11 & Litberg & 65,536 & 27 & Lichtenberg & 244,879\\
12 & Wilsede & 170,876 & 28 & Brocken & 1150,368\\
13 & Bottel & 52,441 & 29 & Hils & 429,082\\
14 & Bullerberg & 53,310 & 30 & Hohehagen & 504,260\\
15 & Hamburg & 129,370 & 31 & Meridianzeichen & 221,112\\
16 & Elmhorst & 89,990 & 32 & Göttingen & 155,836\\[4pt]
\end{tabular}\par}
\par\smallskip\noindent\hspace*{.03\linewidth}{\footnotesize IX.}\hfill{\footnotesize 57}\hspace*{.12\linewidth}\par

\par\clearpage
\noindent\makebox[0pt][l]{\fontsize{13}{16}\selectfont 450}\hfill{\small\textbf{POSTHUMOUS PAPERS.}}\hfill\null\par\medskip
\par\noindent\hspace*{\parindent}[Corrections must now be added to the assumed values for the heights.
They are denoted by \((i)\), where \(i\) is the number of the triangulation point.
Using the height corrections along a chain of lines
running from Langwarden to Göttingen, these corrections can be expressed in terms of new corrections
\(a,\ b,\ c\ .\ .\ .\). For example, the correction for triangulation point
\(20=(20)=(12)+t=(8)+l+t=(6)+g+l+t\)
\par\noindent\hfill \(=(4)+e+g+l+t=(1)+c+e+g+l+t.\)]
\par\medskip
{\centering
\setlength{\tabcolsep}{.57em}
\begin{tabular}{|r|l||r|l||r|l||r|l|}
\noalign{\hrule height .8pt}
\multicolumn{1}{|c|}{\shortstack{\small Triangle\\\small point}}&
\multicolumn{1}{c||}{\shortstack{\rule{0pt}{1pt}\\\small Corr. to\\\small assumed\\\small height}}&
\multicolumn{1}{c|}{\shortstack{\small Triangle\\\small point}}&
\multicolumn{1}{c||}{\shortstack{\rule{0pt}{1pt}\\\small Corr. to\\\small assumed\\\small height}}&
\multicolumn{1}{c|}{\shortstack{\small Triangle\\\small point}}&
\multicolumn{1}{c||}{\shortstack{\rule{0pt}{1pt}\\\small Corr. to\\\small assumed\\\small height}}&
\multicolumn{1}{c|}{\shortstack{\small Triangle\\\small point}}&
\multicolumn{1}{c|}{\shortstack{\rule{0pt}{1pt}\\\small Corr. to\\\small assumed\\\small height}}\\[3pt]
\hline
1 & \((1)\) & 9 & \((8)+h\) & 17 & \((12)+q\) & 25 & \((23)+y\)\\
2 & \((1)+a\) & 10 & \((8)+i\) & 18 & \((12)+r\) & 26 & \((23)+z\)\\
3 & \((1)+b\) & 11 & \((8)+k\) & 19 & \((12)+s\) & 27 & \((23)+\alpha\)\\
4 & \((1)+c\) & 12 & \((8)+l\) & 20 & \((12)+t\) & 28 & \((27)+\beta\)\\
5 & \((4)+d\) & 13 & \((12)+m\) & 21 & \((12)+u\) & 29 & \((27)+\gamma\)\\
6 & \((4)+e\) & 14 & \((12)+n\) & 22 & \((12)+v\) & 30 & \((29)+\delta\)\\
7 & \((6)+f\) & 15 & \((12)+o\) & 23 & \((12)+w\) & 31 & \((29)+\varepsilon\)\\
8 & \((6)+g\) & 16 & \((12)+p\) & 24 & \((23)+x\) & 32 & \((31)+\zeta\)\\
\end{tabular}\par}

\par\medskip
\par\noindent\hspace*{\parindent}[The following values for the height differences between the endpoints
of the triangle sides were derived from the height measurements themselves.]


\par\medskip
{\centering
\setlength{\tabcolsep}{.24em}
\begin{tabular}{|c|r||c|r||c|r||c|r|}
\noalign{\hrule height .8pt}
\multicolumn{1}{|c|}{\shortstack{\small Triangle\\\small side}}&\multicolumn{1}{c||}{\shortstack{\rule{0pt}{1pt}\\\small Observed\\\small height diff.\\[2pt]\small (metres)}}&
\multicolumn{1}{c|}{\shortstack{\small Triangle\\\small side}}&\multicolumn{1}{c||}{\shortstack{\rule{0pt}{1pt}\\\small Observed\\\small height diff.\\[2pt]\small (metres)}}&
\multicolumn{1}{c|}{\shortstack{\small Triangle\\\small side}}&\multicolumn{1}{c||}{\shortstack{\rule{0pt}{1pt}\\\small Observed\\\small height diff.\\[2pt]\small (metres)}}&
\multicolumn{1}{c|}{\shortstack{\small Triangle\\\small side}}&\multicolumn{1}{c|}{\shortstack{\rule{0pt}{1pt}\\\small Observed\\\small height diff.\\[2pt]\small (metres)}}\\[2pt]
\hline
\(1\ .\ 2\) & \(+7{,}165\) & \(7\ .\ 8\) & \(-40{,}175\) & \(10\ .\ 14\) & \(+2{,}764\) & \(12\ .\ 21\) & \(-49{,}371\)\\
\(1\ .\ 3\) & \(+15{,}460\) & \(7\ .\ 9\) & \(-9{,}680\) & \(11\ .\ 12\) & \(+105{,}997\) & \(12\ .\ 22\) & \(-50{,}216\)\\
\(1\ .\ 4\) & \(+1{,}147\) & \(7\ .\ 10\) & \(-32{,}720\) & \(11\ .\ 15\) & \(+64{,}170\) & \(12\ .\ 23\) & \(-20{,}216\)\\
\(2\ .\ 3\) & \(+3{,}829\) & \(8\ .\ 9\) & \(+28{,}309\) & \(11\ .\ 16\) & \(+23{,}830\) & \(13\ .\ 14\) & \(+0{,}767\)\\
\(3\ .\ 4\) & \(-13{,}764\) & \(8\ .\ 10\) & \(+3{,}012\) & \(12\ .\ 13\) & \(-117{,}063\) & \(14\ .\ 16\) & \(+35{,}994\)\\
\(3\ .\ 5\) & \(+12{,}441\) & \(8\ .\ 11\) & \(+19{,}950\) & \(12\ .\ 14\) & \(-117{,}233\) & \(15\ .\ 18\) & \(-13{,}614\)\\
\(4\ .\ 5\) & \(+26{,}921\) & \(8\ .\ 12\) & \(+126{,}432\) & \(12\ .\ 15\) & \(-41{,}959\) & \(15\ .\ 19\) & \(-11{,}395\)\\
\(4\ .\ 6\) & \(+16{,}112\) & \(9\ .\ 12\) & \(+102{,}028\) & \(12\ .\ 16\) & \(-80{,}116\) & \(16\ .\ 23\) & \(+60{,}280\)\\
\(5\ .\ 6\) & \(-4{,}760\) & \(9\ .\ 13\) & \(-21{,}725\) & \(12\ .\ 17\) & \(-93{,}862\) & \(17\ .\ 18\) & \(+40{,}056\)\\
\(5\ .\ 7\) & \(+33{,}760\) & \(10\ .\ 11\) & \(+17{,}039\) & \(12\ .\ 18\) & \(-54{,}149\) & \(18\ .\ 19\) & \(-0{,}320\)\\
\(6\ .\ 7\) & \(+36{,}112\) & \(10\ .\ 12\) & \(+119{,}407\) & \(12\ .\ 19\) & \(-54{,}110\) & \(19\ .\ 20\) & \(-12{,}925\)\\
\(6\ .\ 8\) & \(+0{,}953\) & \(10\ .\ 13\) & \(+3{,}163\) & \(12\ .\ 20\) & \(-66{,}611\) & \(20\ .\ 21\) & \(+16{,}157\)\\
\end{tabular}\par}

\par\clearpage
\noindent\null\hfill{\small\textbf{HEIGHT MEASUREMENTS.}}\hfill\makebox[0pt][r]{\fontsize{13}{16}\selectfont 451}\par\medskip

\par\medskip
{\centering
\setlength{\tabcolsep}{.24em}
\begin{tabular}{|c|r||c|r||c|r||c|r|}
\noalign{\hrule height .8pt}
\multicolumn{1}{|c|}{\shortstack{\small Triangle\\\small side}}&\multicolumn{1}{c||}{\shortstack{\rule{0pt}{1pt}\\\small Observed\\\small height diff.\\[2pt]\small (metres)}}&
\multicolumn{1}{c|}{\shortstack{\small Triangle\\\small side}}&\multicolumn{1}{c||}{\shortstack{\rule{0pt}{1pt}\\\small Observed\\\small height diff.\\[2pt]\small (metres)}}&
\multicolumn{1}{c|}{\shortstack{\small Triangle\\\small side}}&\multicolumn{1}{c||}{\shortstack{\rule{0pt}{1pt}\\\small Observed\\\small height diff.\\[2pt]\small (metres)}}&
\multicolumn{1}{c|}{\shortstack{\small Triangle\\\small side}}&\multicolumn{1}{c|}{\shortstack{\rule{0pt}{1pt}\\\small Observed\\\small height diff.\\[2pt]\small (metres)}}\\[2pt]
\hline
\(20\ .\ 23\) & \(+45{,}529\) & \(23\ .\ 25\) & \(-73{,}042\) & \(26\ .\ 27\) & \(-64{,}712\) & \(29\ .\ 30\) & \(+76{,}108\)\\
\(21\ .\ 22\) & \(-0{,}489\) & \(23\ .\ 26\) & \(+156{,}588\) & \(26\ .\ 29\) & \(+120{,}811\) & \(29\ .\ 31\) & \(-206{,}580\)\\
\(21\ .\ 23\) & \(+32{,}189\) & \(23\ .\ 27\) & \(+94{,}387\) & \(27\ .\ 28\) & \(+904{,}843\) & \(30\ .\ 31\) & \(-283{,}596\)\\
\(22\ .\ 23\) & \(+28{,}984\) & \(24\ .\ 25\) & \(-17{,}355\) & \(27\ .\ 29\) & \(+185{,}210\) & \(30\ .\ 32\) & \(-349{,}366\)\\
\(22\ .\ 24\) & \(-25{,}602\) & \(25\ .\ 26\) & \(+229{,}312\) & \(28\ .\ 29\) & \(-719{,}612\) & \(31\ .\ 32\) & \(-64{,}334\)\\
\(23\ .\ 24\) & \(-58{,}265\) & \(25\ .\ 27\) & \(+169{,}527\) & \(28\ .\ 30\) & \(-648{,}427\) &  & \\
\end{tabular}\par}
\par\medskip
\par\noindent\hspace*{\parindent}[The error equations can now be obtained using these 3 tables.
For example, the height observation along side \(1\ .\ 2\), with \(v_{1.2}\)
denoting the observational error, gives
\[
\begin{aligned}
+7{,}165+v_{1.2}&=+33{,}419+(2)-24{,}726-(1)\\
v_{1.2}&=+1{,}528+a,
\end{aligned}
\]
\noindent and along side \(11\ .\ 15\):
\[
\begin{aligned}
+64{,}170+v_{11.15}&=+63{,}834+(12)+o-(8)-k\\
v_{11.15}&=-0{,}336-k+l+o.
\end{aligned}
\]
\par\noindent\hspace*{\parindent}The error equations are:]
\[
\begin{aligned}
\makebox[0pt][r]{[}v_{1.2}&=+1{,}528+a &\qquad \makebox[0pt][r]{[}v_{7.8}&=+2{,}113-f+g\\
v_{1.3}&=-1{,}410+b &\qquad v_{7.9}&=-0{,}866-f+g+h\\
v_{1.4}&=-0{,}117+c &\qquad v_{7.10}&=-0{,}531-f+g+i\\
v_{2.3}&=+1{,}528-a+b &\qquad v_{8.9}&=-0{,}793+h\\
v_{3.4}&=+0{,}744-b+c &\qquad v_{8.10}&=+1{,}799+i\\
v_{3.5}&=-0{,}627-b+c+d &\qquad v_{8.11}&=+0{,}766+k\\
v_{4.5}&=-2{,}087+d &\qquad v_{8.12}&=-0{,}376+l\\
v_{4.6}&=+2{,}716+e &\qquad v_{9.12}&=-3{,}488-h+l\\
v_{5.6}&=-1{,}246-d+e &\qquad v_{9.13}&=+1{,}830-h+l+m\\
v_{5.7}&=-1{,}468-d+e+f &\qquad v_{10.11}&=-1{,}134-i+k\\
v_{6.7}&=+2{,}186+f &\qquad v_{10.12}&=+1{,}838-i+l\\
v_{6.8}&=]-0{,}717+g &\qquad v_{10.13}&=]-0{,}353-i+l+m\\
\end{aligned}
\]
\par\smallskip\noindent\hfill{\footnotesize 57\textsuperscript{*}}\hspace*{.12\linewidth}\par

\par\clearpage
\noindent\makebox[0pt][l]{\fontsize{13}{16}\selectfont 452}\hfill{\small\textbf{POSTHUMOUS PAPERS.}}\hfill\null\par\medskip
\[
\begin{aligned}
\makebox[0pt][r]{[}v_{10.14}&=+0{,}915-i+l+n &\qquad \makebox[0pt][r]{[}v_{20.23}&=+0{,}766-t+w\\
v_{11.12}&=-0{,}657-k+l &\qquad v_{21.22}&=+0{,}545-u+v\\
v_{11.15}&=-0{,}336-k+l+o &\qquad v_{21.23}&=-1{,}827-u+w\\
v_{11.16}&=+0{,}624-k+l+p &\qquad v_{22.23}&=+1{,}322-v+w\\
v_{12.13}&=-1{,}372+m &\qquad v_{22.24}&=-0{,}940-v+w+x\\
v_{12.14}&=-0{,}333+n &\qquad v_{23.24}&=+1{,}417+x\\
v_{12.15}&=+0{,}453+o &\qquad v_{23.25}&=-0{,}684+y\\
v_{12.16}&=-0{,}770+p &\qquad v_{23.26}&=+0{,}519+z\\
v_{12.17}&=-0{,}182+q &\qquad v_{23.27}&=-0{,}313+\alpha\\
v_{12.18}&=-0{,}020+r &\qquad v_{24.25}&=+0{,}477-x+y\\
v_{12.19}&=+0{,}374+s &\qquad v_{25.26}&=+1{,}521-y+z\\
v_{12.20}&=+0{,}245+t &\qquad v_{25.27}&=-1{,}727-y+\alpha\\
v_{12.21}&=-1{,}062+u &\qquad v_{26.27}&=+1{,}679-z+\alpha\\
v_{12.22}&=-0{,}161+v &\qquad v_{26.29}&=+0{,}359-z+\alpha+\gamma\\
v_{12.23}&=+0{,}145+w &\qquad v_{27.28}&=+0{,}646+\beta\\
v_{13.14}&=+0{,}102-m+n &\qquad v_{27.29}&=-1{,}007+\gamma\\
v_{14.16}&=+0{,}686-n+p &\qquad v_{28.29}&=-1{,}674-\beta+\gamma\\
v_{15.18}&=+0{,}951-o+r &\qquad v_{28.30}&=+2{,}319-\beta+\gamma+\delta\\
v_{15\;19}&=-0{,}835-o+s &\qquad v_{29.30}&=-0{,}930+\delta\\
v_{16.23}&=+0{,}535-p+w &\qquad v_{29.31}&=-1{,}390+\varepsilon\\
v_{17.18}&=-0{,}181-q+r &\qquad v_{30.31}&=+0{,}448-\delta+\varepsilon\\
v_{18.19}&=+0{,}753-r+s &\qquad v_{30.32}&=+0{,}942-\delta+\varepsilon+\zeta\\
v_{19.20}&=+0{,}295-s+t &\qquad v_{31.32}&=]-0{,}942+\zeta.\\
v_{20.21}&=]-0{,}224-t+u &\qquad &\\
\end{aligned}
\]
\par\medskip
\par\noindent\hspace*{\parindent}[This gives, if the weights of the error equations are all
taken to be 1, the following normal equations, in which
the constants are expressed in units of the third decimal place.]

\par\clearpage
\noindent\null\hfill{\small\textbf{HEIGHT MEASUREMENTS.}}\hfill\makebox[0pt][r]{\fontsize{13}{16}\selectfont 453}\par\medskip
{\setlength{\tabcolsep}{0pt}
\renewcommand{\arraystretch}{1.25}
\newcommand{\GaussIXNormalTerm}[1]{\makebox[1.95em][r]{\(#1\)}}
\par\noindent
\begin{tabular}{@{}r*{14}{r}@{}}
\([0=\) & \GaussIXNormalTerm{\ast} & \GaussIXNormalTerm{+2a} & \GaussIXNormalTerm{-b} & \GaussIXNormalTerm{\ast} & \GaussIXNormalTerm{\ast} & \GaussIXNormalTerm{\ast} & \GaussIXNormalTerm{\ast} & \GaussIXNormalTerm{\ast} & \GaussIXNormalTerm{\ast} & \GaussIXNormalTerm{\ast} & \GaussIXNormalTerm{\ast} & \GaussIXNormalTerm{\ast} & \GaussIXNormalTerm{\ast} & \GaussIXNormalTerm{\ast}\\
\(0=\) & \GaussIXNormalTerm{+1} & \GaussIXNormalTerm{-a} & \GaussIXNormalTerm{+4b} & \GaussIXNormalTerm{-2c} & \GaussIXNormalTerm{-d} & \GaussIXNormalTerm{\ast} & \GaussIXNormalTerm{\ast} & \GaussIXNormalTerm{\ast} & \GaussIXNormalTerm{\ast} & \GaussIXNormalTerm{\ast} & \GaussIXNormalTerm{\ast} & \GaussIXNormalTerm{\ast} & \GaussIXNormalTerm{\ast} & \GaussIXNormalTerm{\ast}\\
\(0=\) & \GaussIXNormalTerm{\ast} & \GaussIXNormalTerm{\ast} & \GaussIXNormalTerm{-2b} & \GaussIXNormalTerm{+3c} & \GaussIXNormalTerm{+d} & \GaussIXNormalTerm{\ast} & \GaussIXNormalTerm{\ast} & \GaussIXNormalTerm{\ast} & \GaussIXNormalTerm{\ast} & \GaussIXNormalTerm{\ast} & \GaussIXNormalTerm{\ast} & \GaussIXNormalTerm{\ast} & \GaussIXNormalTerm{\ast} & \GaussIXNormalTerm{\ast}\\
\(0=\) & \GaussIXNormalTerm{\ast} & \GaussIXNormalTerm{\ast} & \GaussIXNormalTerm{-b} & \GaussIXNormalTerm{+c} & \GaussIXNormalTerm{+4d} & \GaussIXNormalTerm{-2e} & \GaussIXNormalTerm{-f} & \GaussIXNormalTerm{\ast} & \GaussIXNormalTerm{\ast} & \GaussIXNormalTerm{\ast} & \GaussIXNormalTerm{\ast} & \GaussIXNormalTerm{\ast} & \GaussIXNormalTerm{\ast} & \GaussIXNormalTerm{\ast}\\
\(0=\) & \GaussIXNormalTerm{+2} & \GaussIXNormalTerm{\ast} & \GaussIXNormalTerm{\ast} & \GaussIXNormalTerm{\ast} & \GaussIXNormalTerm{-2d} & \GaussIXNormalTerm{+3e} & \GaussIXNormalTerm{+f} & \GaussIXNormalTerm{\ast} & \GaussIXNormalTerm{\ast} & \GaussIXNormalTerm{\ast} & \GaussIXNormalTerm{\ast} & \GaussIXNormalTerm{\ast} & \GaussIXNormalTerm{\ast} & \GaussIXNormalTerm{\ast}\\
\(0=\) & \GaussIXNormalTerm{+2} & \GaussIXNormalTerm{\ast} & \GaussIXNormalTerm{\ast} & \GaussIXNormalTerm{\ast} & \GaussIXNormalTerm{-d} & \GaussIXNormalTerm{+e} & \GaussIXNormalTerm{+5f} & \GaussIXNormalTerm{-3g} & \GaussIXNormalTerm{-h} & \GaussIXNormalTerm{-i} & \GaussIXNormalTerm{\ast} & \GaussIXNormalTerm{\ast} & \GaussIXNormalTerm{\ast} & \GaussIXNormalTerm{\ast}\\
\(0=\) & \GaussIXNormalTerm{-1} & \GaussIXNormalTerm{\ast} & \GaussIXNormalTerm{\ast} & \GaussIXNormalTerm{\ast} & \GaussIXNormalTerm{\ast} & \GaussIXNormalTerm{\ast} & \GaussIXNormalTerm{-3f} & \GaussIXNormalTerm{+4g} & \GaussIXNormalTerm{+h} & \GaussIXNormalTerm{+i} & \GaussIXNormalTerm{\ast} & \GaussIXNormalTerm{\ast} & \GaussIXNormalTerm{\ast} & \GaussIXNormalTerm{\ast}\\
\(0=\) & \GaussIXNormalTerm{-1} & \GaussIXNormalTerm{\ast} & \GaussIXNormalTerm{\ast} & \GaussIXNormalTerm{\ast} & \GaussIXNormalTerm{\ast} & \GaussIXNormalTerm{\ast} & \GaussIXNormalTerm{-f} & \GaussIXNormalTerm{+g} & \GaussIXNormalTerm{+4h} & \GaussIXNormalTerm{\ast} & \GaussIXNormalTerm{\ast} & \GaussIXNormalTerm{-2l} & \GaussIXNormalTerm{-m} & \GaussIXNormalTerm{\ast}\\
\(0=\) & \GaussIXNormalTerm{+2} & \GaussIXNormalTerm{\ast} & \GaussIXNormalTerm{\ast} & \GaussIXNormalTerm{\ast} & \GaussIXNormalTerm{\ast} & \GaussIXNormalTerm{\ast} & \GaussIXNormalTerm{-f} & \GaussIXNormalTerm{+g} & \GaussIXNormalTerm{\ast} & \GaussIXNormalTerm{+6i} & \GaussIXNormalTerm{-k} & \GaussIXNormalTerm{-3l} & \GaussIXNormalTerm{-m} & \GaussIXNormalTerm{-n}\\
\end{tabular}\par\smallskip
\par\noindent
\begin{tabular}{@{}r*{16}{r}@{}}
\(0=\) & \GaussIXNormalTerm{+1} & \GaussIXNormalTerm{\ast} & \GaussIXNormalTerm{-i} & \GaussIXNormalTerm{+5k} & \GaussIXNormalTerm{-3l} & \GaussIXNormalTerm{\ast} & \GaussIXNormalTerm{\ast} & \GaussIXNormalTerm{-o} & \GaussIXNormalTerm{-p} & \GaussIXNormalTerm{\ast} & \GaussIXNormalTerm{\ast} & \GaussIXNormalTerm{\ast} & \GaussIXNormalTerm{\ast} & \GaussIXNormalTerm{\ast} & \GaussIXNormalTerm{\ast} & \GaussIXNormalTerm{\ast}\\
\(0=\) & \GaussIXNormalTerm{-3} & \GaussIXNormalTerm{-2h} & \GaussIXNormalTerm{-3i} & \GaussIXNormalTerm{-3k} & \GaussIXNormalTerm{+9l} & \GaussIXNormalTerm{+2m} & \GaussIXNormalTerm{+n} & \GaussIXNormalTerm{+o} & \GaussIXNormalTerm{+p} & \GaussIXNormalTerm{\ast} & \GaussIXNormalTerm{\ast} & \GaussIXNormalTerm{\ast} & \GaussIXNormalTerm{\ast} & \GaussIXNormalTerm{\ast} & \GaussIXNormalTerm{\ast} & \GaussIXNormalTerm{\ast}\\
\(0=\) & \GaussIXNormalTerm{+3} & \GaussIXNormalTerm{-h} & \GaussIXNormalTerm{-i} & \GaussIXNormalTerm{\ast} & \GaussIXNormalTerm{+2l} & \GaussIXNormalTerm{+4m} & \GaussIXNormalTerm{-n} & \GaussIXNormalTerm{\ast} & \GaussIXNormalTerm{\ast} & \GaussIXNormalTerm{\ast} & \GaussIXNormalTerm{\ast} & \GaussIXNormalTerm{\ast} & \GaussIXNormalTerm{\ast} & \GaussIXNormalTerm{\ast} & \GaussIXNormalTerm{\ast} & \GaussIXNormalTerm{\ast}\\
\(0=\) & \GaussIXNormalTerm{-2} & \GaussIXNormalTerm{\ast} & \GaussIXNormalTerm{-i} & \GaussIXNormalTerm{\ast} & \GaussIXNormalTerm{+l} & \GaussIXNormalTerm{-m} & \GaussIXNormalTerm{+4n} & \GaussIXNormalTerm{\ast} & \GaussIXNormalTerm{-p} & \GaussIXNormalTerm{\ast} & \GaussIXNormalTerm{\ast} & \GaussIXNormalTerm{\ast} & \GaussIXNormalTerm{\ast} & \GaussIXNormalTerm{\ast} & \GaussIXNormalTerm{\ast} & \GaussIXNormalTerm{\ast}\\
\(0=\) & \GaussIXNormalTerm{+1} & \GaussIXNormalTerm{\ast} & \GaussIXNormalTerm{\ast} & \GaussIXNormalTerm{-k} & \GaussIXNormalTerm{+l} & \GaussIXNormalTerm{\ast} & \GaussIXNormalTerm{\ast} & \GaussIXNormalTerm{+4o} & \GaussIXNormalTerm{\ast} & \GaussIXNormalTerm{\ast} & \GaussIXNormalTerm{-r} & \GaussIXNormalTerm{-s} & \GaussIXNormalTerm{\ast} & \GaussIXNormalTerm{\ast} & \GaussIXNormalTerm{\ast} & \GaussIXNormalTerm{\ast}\\
\(0=\) & \GaussIXNormalTerm{+5} & \GaussIXNormalTerm{\ast} & \GaussIXNormalTerm{\ast} & \GaussIXNormalTerm{-k} & \GaussIXNormalTerm{+l} & \GaussIXNormalTerm{\ast} & \GaussIXNormalTerm{-n} & \GaussIXNormalTerm{\ast} & \GaussIXNormalTerm{+4p} & \GaussIXNormalTerm{\ast} & \GaussIXNormalTerm{\ast} & \GaussIXNormalTerm{\ast} & \GaussIXNormalTerm{\ast} & \GaussIXNormalTerm{\ast} & \GaussIXNormalTerm{\ast} & \GaussIXNormalTerm{-w}\\
\(0=\) & \GaussIXNormalTerm{-1} & \GaussIXNormalTerm{\ast} & \GaussIXNormalTerm{\ast} & \GaussIXNormalTerm{\ast} & \GaussIXNormalTerm{\ast} & \GaussIXNormalTerm{\ast} & \GaussIXNormalTerm{\ast} & \GaussIXNormalTerm{\ast} & \GaussIXNormalTerm{\ast} & \GaussIXNormalTerm{+2q} & \GaussIXNormalTerm{-r} & \GaussIXNormalTerm{\ast} & \GaussIXNormalTerm{\ast} & \GaussIXNormalTerm{\ast} & \GaussIXNormalTerm{\ast} & \GaussIXNormalTerm{\ast}\\
\(0=\) & \GaussIXNormalTerm{-3} & \GaussIXNormalTerm{\ast} & \GaussIXNormalTerm{\ast} & \GaussIXNormalTerm{\ast} & \GaussIXNormalTerm{\ast} & \GaussIXNormalTerm{\ast} & \GaussIXNormalTerm{\ast} & \GaussIXNormalTerm{-o} & \GaussIXNormalTerm{\ast} & \GaussIXNormalTerm{-q} & \GaussIXNormalTerm{+4r} & \GaussIXNormalTerm{-s} & \GaussIXNormalTerm{\ast} & \GaussIXNormalTerm{\ast} & \GaussIXNormalTerm{\ast} & \GaussIXNormalTerm{\ast}\\
\(0=\) & \GaussIXNormalTerm{-3} & \GaussIXNormalTerm{\ast} & \GaussIXNormalTerm{\ast} & \GaussIXNormalTerm{\ast} & \GaussIXNormalTerm{\ast} & \GaussIXNormalTerm{\ast} & \GaussIXNormalTerm{\ast} & \GaussIXNormalTerm{-o} & \GaussIXNormalTerm{\ast} & \GaussIXNormalTerm{\ast} & \GaussIXNormalTerm{-r} & \GaussIXNormalTerm{+4s} & \GaussIXNormalTerm{-t} & \GaussIXNormalTerm{\ast} & \GaussIXNormalTerm{\ast} & \GaussIXNormalTerm{\ast}\\
\(0=\) & \GaussIXNormalTerm{-2} & \GaussIXNormalTerm{\ast} & \GaussIXNormalTerm{\ast} & \GaussIXNormalTerm{\ast} & \GaussIXNormalTerm{\ast} & \GaussIXNormalTerm{\ast} & \GaussIXNormalTerm{\ast} & \GaussIXNormalTerm{\ast} & \GaussIXNormalTerm{\ast} & \GaussIXNormalTerm{\ast} & \GaussIXNormalTerm{\ast} & \GaussIXNormalTerm{-s} & \GaussIXNormalTerm{+4t} & \GaussIXNormalTerm{-u} & \GaussIXNormalTerm{\ast} & \GaussIXNormalTerm{-w}\\
\(0=\) & \GaussIXNormalTerm{-4} & \GaussIXNormalTerm{\ast} & \GaussIXNormalTerm{\ast} & \GaussIXNormalTerm{\ast} & \GaussIXNormalTerm{\ast} & \GaussIXNormalTerm{\ast} & \GaussIXNormalTerm{\ast} & \GaussIXNormalTerm{\ast} & \GaussIXNormalTerm{\ast} & \GaussIXNormalTerm{\ast} & \GaussIXNormalTerm{\ast} & \GaussIXNormalTerm{\ast} & \GaussIXNormalTerm{-t} & \GaussIXNormalTerm{+4u} & \GaussIXNormalTerm{-v} & \GaussIXNormalTerm{-w}\\
\end{tabular}\par\smallskip
\par\noindent
\begin{tabular}{@{}r*{13}{r}@{}}
\(0=\) & \GaussIXNormalTerm{+2} & \GaussIXNormalTerm{\ast} & \GaussIXNormalTerm{\ast} & \GaussIXNormalTerm{\ast} & \GaussIXNormalTerm{\ast} & \GaussIXNormalTerm{\ast} & \GaussIXNormalTerm{-u} & \GaussIXNormalTerm{+4v} & \GaussIXNormalTerm{-2w} & \GaussIXNormalTerm{-x} & \GaussIXNormalTerm{\ast} & \GaussIXNormalTerm{\ast} & \GaussIXNormalTerm{\ast}\\
\(0=\) & \GaussIXNormalTerm{+1} & \GaussIXNormalTerm{-p} & \GaussIXNormalTerm{\ast} & \GaussIXNormalTerm{\ast} & \GaussIXNormalTerm{\ast} & \GaussIXNormalTerm{-t} & \GaussIXNormalTerm{-u} & \GaussIXNormalTerm{-2v} & \GaussIXNormalTerm{+6w} & \GaussIXNormalTerm{+x} & \GaussIXNormalTerm{\ast} & \GaussIXNormalTerm{\ast} & \GaussIXNormalTerm{\ast}\\
\(0=\) & \GaussIXNormalTerm{\ast} & \GaussIXNormalTerm{\ast} & \GaussIXNormalTerm{\ast} & \GaussIXNormalTerm{\ast} & \GaussIXNormalTerm{\ast} & \GaussIXNormalTerm{\ast} & \GaussIXNormalTerm{\ast} & \GaussIXNormalTerm{-v} & \GaussIXNormalTerm{+w} & \GaussIXNormalTerm{+3x} & \GaussIXNormalTerm{-y} & \GaussIXNormalTerm{\ast} & \GaussIXNormalTerm{\ast}\\
\(0=\) & \GaussIXNormalTerm{-1} & \GaussIXNormalTerm{\ast} & \GaussIXNormalTerm{\ast} & \GaussIXNormalTerm{\ast} & \GaussIXNormalTerm{\ast} & \GaussIXNormalTerm{\ast} & \GaussIXNormalTerm{\ast} & \GaussIXNormalTerm{\ast} & \GaussIXNormalTerm{\ast} & \GaussIXNormalTerm{-x} & \GaussIXNormalTerm{+4y} & \GaussIXNormalTerm{-z} & \GaussIXNormalTerm{-\alpha}\\
\end{tabular}\par\smallskip
\par\noindent
\begin{tabular}{@{}r*{9}{r}@{}}
\(0=\) & \GaussIXNormalTerm{+2} & \GaussIXNormalTerm{-y} & \GaussIXNormalTerm{+4z} & \GaussIXNormalTerm{-2\alpha} & \GaussIXNormalTerm{\ast} & \GaussIXNormalTerm{-\gamma} & \GaussIXNormalTerm{\ast} & \GaussIXNormalTerm{\ast} & \GaussIXNormalTerm{\ast}\\
\(0=\) & \GaussIXNormalTerm{-2} & \GaussIXNormalTerm{-y} & \GaussIXNormalTerm{-2z} & \GaussIXNormalTerm{+4\alpha} & \GaussIXNormalTerm{\ast} & \GaussIXNormalTerm{+\gamma} & \GaussIXNormalTerm{\ast} & \GaussIXNormalTerm{\ast} & \GaussIXNormalTerm{\ast}\\
\(0=\) & \GaussIXNormalTerm{+1} & \GaussIXNormalTerm{\ast} & \GaussIXNormalTerm{\ast} & \GaussIXNormalTerm{\ast} & \GaussIXNormalTerm{+3\beta} & \GaussIXNormalTerm{-2\gamma} & \GaussIXNormalTerm{-\delta} & \GaussIXNormalTerm{\ast} & \GaussIXNormalTerm{\ast}\\
\(0=\) & \GaussIXNormalTerm{-3} & \GaussIXNormalTerm{\ast} & \GaussIXNormalTerm{-z} & \GaussIXNormalTerm{+\alpha} & \GaussIXNormalTerm{-2\beta} & \GaussIXNormalTerm{+4\gamma} & \GaussIXNormalTerm{+\delta} & \GaussIXNormalTerm{\ast} & \GaussIXNormalTerm{\ast}\\
\(0=\) & \GaussIXNormalTerm{-1} & \GaussIXNormalTerm{\ast} & \GaussIXNormalTerm{\ast} & \GaussIXNormalTerm{\ast} & \GaussIXNormalTerm{-\beta} & \GaussIXNormalTerm{+\gamma} & \GaussIXNormalTerm{+4\delta} & \GaussIXNormalTerm{-2\varepsilon} & \GaussIXNormalTerm{-\zeta}\\
\(0=\) & \GaussIXNormalTerm{\ast} & \GaussIXNormalTerm{\ast} & \GaussIXNormalTerm{\ast} & \GaussIXNormalTerm{\ast} & \GaussIXNormalTerm{\ast} & \GaussIXNormalTerm{\ast} & \GaussIXNormalTerm{-2\delta} & \GaussIXNormalTerm{+3\varepsilon} & \GaussIXNormalTerm{+\zeta}\\
\(0=\)\makebox[0pt][l]{]} & \GaussIXNormalTerm{\ast} & \GaussIXNormalTerm{\ast} & \GaussIXNormalTerm{\ast} & \GaussIXNormalTerm{\ast} & \GaussIXNormalTerm{\ast} & \GaussIXNormalTerm{\ast} & \GaussIXNormalTerm{-\delta} & \GaussIXNormalTerm{+\varepsilon} & \GaussIXNormalTerm{+2\zeta.}\\
\end{tabular}\par
}

\par\clearpage
\noindent\makebox[0pt][l]{\fontsize{13}{16}\selectfont 454}\hfill{\small\textbf{POSTHUMOUS PAPERS.\quad HEIGHT MEASUREMENTS.}}\hfill\null\par\medskip
\par\noindent\hspace*{\parindent}The following corrections can still be applied:
\[
\begin{aligned}
f&=-0{,}001 & l&=-0{,}001 & q&=+0{,}001 & t&=+0{,}001\\
i&=-0{,}001 & m&=-0{,}001 & r&=+0{,}001 & u&=+0{,}001\\
k&=-0{,}001 & p&=-0{,}001 & s&=+0{,}001 & [\gamma&=+0{,}001].
\end{aligned}
\]
\par\noindent\hspace*{\parindent}[Accordingly, the following corrections must still be added to the assumed heights of the
triangulation points (Table I), if at the same time
\((1)=0\) is set:
\[
\begin{aligned}
(1)&=\phantom{-}0 & (9)&=\phantom{-}0 & (17)&=\phantom{-}0 & (25)&=-0{,}001\\
(2)&=\phantom{-}0 & (10)&=-0{,}001 & (18)&=\phantom{-}0 & (26)&=-0{,}001\\
(3)&=\phantom{-}0 & (11)&=-0{,}001 & (19)&=\phantom{-}0 & (27)&=-0{,}001\\
(4)&=\phantom{-}0 & (12)&=-0{,}001 & (20)&=\phantom{-}0 & (28)&=-0{,}001\\
(5)&=\phantom{-}0 & (13)&=-0{,}002 & (21)&=\phantom{-}0 & (29)&=\phantom{-}0\\
(6)&=\phantom{-}0 & (14)&=-0{,}001 & (22)&=-0{,}001 & (30)&=\phantom{-}0\\
(7)&=-0{,}001 & (15)&=-0{,}001 & (23)&=-0{,}001 & (31)&=\phantom{-}0\\
(8)&=\phantom{-}0 & (16)&=-0{,}002 & (24)&=-0{,}001 & (32)&=\phantom{-}0\quad.]
\end{aligned}
\]
\par\medskip
\begin{center}\rule{.18\linewidth}{.4pt}\end{center}
\par\bigskip
\begin{center}\small\textbf{REMARKS.}\end{center}
\begingroup
\footnotesize
\setlength{\abovedisplayskip}{5pt}\setlength{\belowdisplayskip}{5pt}
\par\noindent\hspace*{\parindent}Note [1] is on a separate leaf.
A scribal error was corrected in its formula II) and in the immediately preceding formula for \(h\).
In place of \(\Delta\), the original has \(t\) in both instances.
\par\noindent\hspace*{\parindent}Notes [2] and [4] are taken from 2 notebooks, and Note [3] from a
calculation notebook for the Hanover arc measurement.
\par\noindent\hspace*{\parindent}The formulas in Art. [2] follow directly from the derivations
in Art. [1]. From the formula at the top of p. 446,
\[
\left(\tfrac ab\right)^{1-n}\mathbin{.}\cos\alpha
=\left(\tfrac{a}{a+h}\right)^{1-n}\mathbin{.}\cos\alpha
=\cos(\alpha+(1-n)t)
\]
\noindent it follows for \(\alpha=0\) that:
\[
\left(\tfrac{a}{a+h}\right)^{1-n}=\cos(1-n)t.
\]
\par\noindent\hspace*{\parindent}Writing \(R\) for \(a\), \(v\) for \(t\) and
\(\varepsilon\) for the coefficient of refraction \(n\), one obtains from this,
since \(v-\varepsilon v=\delta\):
\[
\left(\tfrac{R}{R+h}\right)^{1-\varepsilon}=\cos\delta.
\]
\endgroup

\par\clearpage
\noindent\null\hfill{\small\textbf{REMARKS.\quad HEIGHT MEASUREMENTS.}}\hfill\makebox[0pt][r]{\fontsize{13}{16}\selectfont 455}\par\medskip
\begingroup
\footnotesize
\setlength{\abovedisplayskip}{5pt}\setlength{\belowdisplayskip}{5pt}
\par\noindent\hspace*{\parindent}In place of the formula given in Art. [2] for
\(\dfrac{R+h}{R}\), the original had:
\[
\frac{R+h}{R}=
\frac{\cos\dfrac{2-\varepsilon}{2-2\varepsilon}\mathbin{.}\delta}
     {\cos\dfrac{\varepsilon}{2-2\varepsilon}\mathbin{.}\delta}.
\]
\par\noindent\hspace*{\parindent}The formula for \(h'-h\) in Art. [4], which \textsc{Gauss} used to
calculate the height differences between his southern principal triangulation points,
likewise assumes the Earth to be spherical over the length of the light ray
in question. Let \(M\) be its centre and \(R\) its radius,
with \(PM=R+h\) and \(P'M=R+h'\); let \(z\) be the zenith distance of the light ray
from \(P'\) at \(P\) and \(z'\) the zenith distance from \(P\) at \(P'\); let the corresponding
angles of refraction be \(\delta z\) and \(\delta z'\).
\par\noindent\hspace*{\parindent}If one sets
\[
z+\delta z=90^{0}+a\quad\text{and}\quad z'+\delta z'=90^{0}+a',
\]
\noindent then in triangle \(PMP'\):
\[
h'-h=(2R+h+h')\operatorname{tang}\tfrac12(a'-a)\operatorname{tang}\tfrac12(a'+a),
\]
\noindent or, since for small angles \(\operatorname{tang}x=x\,{}^3\!\surd\sec x^2\) approximately,
\[
h'-h=(R+\tfrac12(h+h'))c\mathbin{.}\tfrac12(a'-a)
(\sec\tfrac12(a'-a)\sec\tfrac12(a'+a))^{\frac23},
\]
\noindent where the angle is \(PMP'=c=a'+a\).
\par\noindent\hspace*{\parindent}Because \(Rc=r\), and because for small angles
\[
(\sec\tfrac12(a'-a)\sec\tfrac12(a'+a))^2
=1+\tfrac12(a'a'+aa)=\sec a\sec a'
\]
\noindent also holds, it follows that:
\[
h'-h=\tfrac12(a'-a)(\sec a\sec a')^{\frac13}(r+\tfrac12(h+h')c).
\]
\par\noindent\hspace*{\parindent}\textsc{Gauss} took the refraction angles
\(\delta z\) and \(\delta z'\) to be equal in calculating the height differences.
In most cases, the formula used for this calculation was
\[
h'-h=s\operatorname{tang}\tfrac12(z'-z)
\]
\noindent in which \(s\) denotes the triangle side.
\par\noindent\hspace*{\parindent}The tables in Note [5] on the adjustment of the heights of the
triangulation points were compiled from entries on one page of the notebook entitled:
»Aufsätze, Notizen und Rechnungen, zur Mathematik gehörig«;
the column headings were added. The completion of this
adjustment was announced by \textsc{Gauss} in his letter to \textsc{Olbers} of
2 April 1826 (p. 377). The approximate values used (in the first table)
agree almost exactly with the adjusted values, probably because
an adjustment had already been carried out. Likewise, in the adjustments of the
station observations one frequently finds that \textsc{Gauss}, after completing an
adjustment, set up the error and normal equations once again
(see the station adjustments for Brillit and Wilsede, pp. 265/270).
Several scribal errors have been corrected in the error equations on pp. 451/452.
\par\noindent\hspace*{\parindent}Substituting the values of the corrections into these error equations
gives the mean error of an adjusted height difference
\[
=\surd\frac{100{,}3719}{71-31}=\pm1{,}584\,\text{m}.
\]
\par\noindent\hspace*{\parindent}It should also be mentioned that the angles of elevation were generally obtained from
20 repetitions.
\par\noindent\hspace*{\parindent}The heights reported in the letter to \textsc{Olbers} of 19 February 1825, pp. 373/374,
are preliminary values.\hfill\textsc{Krüger.}\par
\endgroup
\par\medskip
\begin{center}\rule{.24\linewidth}{.4pt}\end{center}

\par\clearpage
\noindent\makebox[0pt][l]{\fontsize{13}{16}\selectfont 456}\hfill{\small\textbf{HEIGHT MEASUREMENTS.}}\hfill\null\par\medskip
\vspace*{.23\textheight}
\begin{center}
[Tables for barometric height measurement.]
\par\bigskip
{\footnotesize
\begin{tabular}{c}
\hline
Astronomisches Jahrbuch für das Jahr 1818.\quad Berlin 1815.\quad pp. 169---172.\\
\hline
\end{tabular}}
\end{center}
\par\noindent\hspace*{\parindent}Measuring heights with a barometer is admittedly no astronomical subject,
yet there will be no reader of the Yearbook to whom it is not also of interest,
and so I believe they will welcome my presenting a small table for this purpose
that I calculated some time ago for my own use. To \textls[120]{me}, it is
more convenient than all the extensive auxiliary tables; it reproduces
\textsc{Laplace}'s expression with complete rigour. On other occasions too, for example in the case of aberration,
nutation and corresponding solar altitudes, the newly devised tables, which in conjunction with
tables of logarithms give the desired result as conveniently as possible, have been well received; I hope
the same will be true of the present tables, which one can easily copy onto
the blank leaf of the table of logarithms one is accustomed to using.
\begin{center}Table I.\end{center}

{\centering
\setlength{\tabcolsep}{.75em}
\begin{tabular}{@{\vrule width .8pt\hspace{\tabcolsep}}r|r@{\hspace{\tabcolsep}\vrule width .8pt\hspace{\tabcolsep}}r|r@{\hspace{\tabcolsep}\vrule width .8pt\hspace{\tabcolsep}}r|r@{\hspace{\tabcolsep}\vrule width .8pt\hspace{\tabcolsep}}r|r@{\hspace{\tabcolsep}\vrule width .8pt}}
\noalign{\hrule height .8pt}
\rule[-.8\baselineskip]{0pt}{2.35\baselineskip}\makebox[2.5em][c]{\(t+t'\)} & \makebox[3.6em][c]{\(A\)} & \rule[-.8\baselineskip]{0pt}{2.35\baselineskip}\makebox[2.5em][c]{\(t+t'\)} & \makebox[3.6em][c]{\(A\)} & \rule[-.8\baselineskip]{0pt}{2.35\baselineskip}\makebox[2.5em][c]{\(t+t'\)} & \makebox[3.6em][c]{\(A\)} & \rule[-.8\baselineskip]{0pt}{2.35\baselineskip}\makebox[2.5em][c]{\(t+t'\)} & \makebox[3.6em][c]{\(A\)}\\
\hline
\rule{0pt}{1.35\baselineskip}\(-10^{0}\) & \(4{,}25337\) & \(-5^{0}\) & \(4{,}25892\) & \(0^{0}\) & \(4{,}26439\) & \(+5^{0}\) & \(4{,}26980\)\\
\(-9\) & \(4{,}25448\) & \(-4\) & \(4{,}26002\) & \(+1\) & \(4{,}26548\) & \(+6\) & \(4{,}27087\)\\
\(-8\) & \(4{,}25560\) & \(-3\) & \(4{,}26111\) & \(+2\) & \(4{,}26658\) & \(+7\) & \(4{,}27195\)\\
\(-7\) & \(4{,}25671\) & \(-2\) & \(4{,}26220\) & \(+3\) & \(4{,}26765\) & \(+8\) & \(4{,}27301\)\\
\(-6\) & \(4{,}25781\) & \(-1\) & \(4{,}26330\) & \(+4\) & \(4{,}26872\) & \(+9\) & \(4{,}27408\)\\
\(-5\) & \(4{,}25892\) & \(0\) & \(4{,}26439\) & \(+5\) & \(4{,}26980\) & \(+10\) & \(4{,}27514\)\\
\end{tabular}\par}

\par\clearpage
\noindent\null\hfill{\small\textbf{HEIGHT MEASUREMENTS.}}\hfill\makebox[0pt][r]{\fontsize{13}{16}\selectfont 457}\par\medskip

{\centering
\setlength{\tabcolsep}{.75em}
\begin{tabular}{@{\vrule width .8pt\hspace{\tabcolsep}}r|r@{\hspace{\tabcolsep}\vrule width .8pt\hspace{\tabcolsep}}r|r@{\hspace{\tabcolsep}\vrule width .8pt\hspace{\tabcolsep}}r|r@{\hspace{\tabcolsep}\vrule width .8pt\hspace{\tabcolsep}}r|r@{\hspace{\tabcolsep}\vrule width .8pt}}
\noalign{\hrule height .8pt}
\rule[-.8\baselineskip]{0pt}{2.35\baselineskip}\makebox[2.5em][c]{\(t+t'\)} & \makebox[3.6em][c]{\(A\)} & \rule[-.8\baselineskip]{0pt}{2.35\baselineskip}\makebox[2.5em][c]{\(t+t'\)} & \makebox[3.6em][c]{\(A\)} & \rule[-.8\baselineskip]{0pt}{2.35\baselineskip}\makebox[2.5em][c]{\(t+t'\)} & \makebox[3.6em][c]{\(A\)} & \rule[-.8\baselineskip]{0pt}{2.35\baselineskip}\makebox[2.5em][c]{\(t+t'\)} & \makebox[3.6em][c]{\(A\)}\\
\hline
\rule{0pt}{1.35\baselineskip}\(+10^{0}\) & \(4{,}27514\) & \(+20^{0}\) & \(4{,}28564\) & \(+30^{0}\) & \(4{,}29588\) & \(+40^{0}\) & \(4{,}30589\)\\
\(11\) & \(4{,}27620\) & \(21\) & \(4{,}28667\) & \(31\) & \(4{,}29689\) & \(41\) & \(4{,}30688\)\\
\(12\) & \(4{,}27726\) & \(22\) & \(4{,}28770\) & \(32\) & \(4{,}29790\) & \(42\) & \(4{,}30787\)\\
\(13\) & \(4{,}27832\) & \(23\) & \(4{,}28874\) & \(33\) & \(4{,}29891\) & \(43\) & \(4{,}30885\)\\
\(14\) & \(4{,}27937\) & \(24\) & \(4{,}28976\) & \(34\) & \(4{,}29991\) & \(44\) & \(4{,}30984\)\\
\(15\) & \(4{,}28042\) & \(25\) & \(4{,}29079\) & \(35\) & \(4{,}30092\) & \(45\) & \(4{,}31082\)\\
\(16\) & \(4{,}28147\) & \(26\) & \(4{,}29181\) & \(36\) & \(4{,}30192\) & \(46\) & \(4{,}31179\)\\
\(17\) & \(4{,}28251\) & \(27\) & \(4{,}29283\) & \(37\) & \(4{,}30291\) & \(47\) & \(4{,}31277\)\\
\(18\) & \(4{,}28356\) & \(28\) & \(4{,}29385\) & \(38\) & \(4{,}30391\) & \(48\) & \(4{,}31374\)\\
\(19\) & \(4{,}28460\) & \(29\) & \(4{,}29487\) & \(39\) & \(4{,}30490\) & \(49\) & \(4{,}31471\)\\
\(20\) & \(4{,}28564\) & \(30\) & \(4{,}29588\) & \(40\) & \(4{,}30589\) & \(50\) & \(4{,}31568\)\\
\end{tabular}\par}
\par\bigskip
\begin{center}
Table II.\par\medskip
Correction to \(A\).\quad Argument: the altitude of the pole.
\end{center}
{\centering
\setlength{\tabcolsep}{.85em}
\begin{tabular}{r|r|r@{\hspace{\tabcolsep}\vrule width .8pt\hspace{\tabcolsep}}r|r|r@{\hspace{\tabcolsep}\vrule width .8pt\hspace{\tabcolsep}}r|r|r}
P. alt. & \makebox[1.7em][c]{\(+\)} & &
P. alt. & \makebox[1.7em][c]{\(+\)} & &
P. alt. & \makebox[1.7em][c]{\(+\)} & \\[4pt]
\(0^{0}\) & \(124\) & \(90^{0}\) & \(15^{0}\) & \(107\) & \(75^{0}\) & \(30^{0}\) & \(62\) & \(60^{0}\)\\
\(1\) & \(123\) & \(89\) & \(16\) & \(105\) & \(74\) & \(31\) & \(58\) & \(59\)\\
\(2\) & \(123\) & \(88\) & \(17\) & \(102\) & \(73\) & \(32\) & \(54\) & \(58\)\\
\(3\) & \(123\) & \(87\) & \(18\) & \(100\) & \(72\) & \(33\) & \(50\) & \(57\)\\
\(4\) & \(122\) & \(86\) & \(19\) & \(.97\) & \(71\) & \(34\) & \(46\) & \(56\)\\
\(5\) & \(122\) & \(85\) & \(20\) & \(95\) & \(70\) & \(35\) & \(42\) & \(55\)\\
\(6\) & \(121\) & \(84\) & \(21\) & \(92\) & \(69\) & \(36\) & \(38\) & \(54\)\\
\(7\) & \(120\) & \(83\) & \(22\) & \(89\) & \(68\) & \(37\) & \(34\) & \(53\)\\
\(8\) & \(119\) & \(82\) & \(23\) & \(86\) & \(67\) & \(38\) & \(30\) & \(52\)\\
\(9\) & \(118\) & \(81\) & \(24\) & \(83\) & \(66\) & \(39\) & \(26\) & \(51\)\\
\(10\) & \(116\) & \(80\) & \(25\) & \(79\) & \(65\) & \(40\) & \(21\) & \(50\)\\
\(11\) & \(115\) & \(79\) & \(26\) & \(76\) & \(64\) & \(41\) & \(17\) & \(49\)\\
\(12\) & \(113\) & \(78\) & \(27\) & \(73\) & \(63\) & \(42\) & \(13\) & \(48\)\\
\(13\) & \(111\) & \(77\) & \(28\) & \(69\) & \(62\) & \(43\) & \(9\) & \(47\)\\
\(14\) & \(109\) & \(76\) & \(29\) & \(65\) & \(61\) & \(44\) & \(4\) & \(46\)\\
\(15\) & \(107\) & \(75\) & \(30\) & \(62\) & \(60\) & \(45\) & \(0\) & \(45\)\\
\rule{0pt}{1.35\baselineskip} & \makebox[1.7em][c]{\(-\)} & P. alt &
& \makebox[1.7em][c]{\(-\)} & P. alt. &
& \makebox[1.7em][c]{\(-\)} & P. alt.\\
\end{tabular}\par}
\par\bigskip\noindent\hspace*{.03\linewidth}{\footnotesize IX.}\hfill{\footnotesize 58}\hspace*{.12\linewidth}\par
